\documentclass[12pt,a4paper]{article}
\usepackage[UTF8]{ctex}
\usepackage{amsmath,amssymb,amsthm}
\usepackage{geometry}

\geometry{margin=2.5cm}

\title{为什么好的运动控制器需要高阻抗（低导纳）？}
\author{第11章补充说明}
\date{}

\begin{document}

\maketitle

\section{问题提出}

在机器人控制中，我们经常说：
\begin{itemize}
\item "一个好的运动控制器的特点是高阻抗（低导纳）"
\item 因为$\Delta X = Y \Delta F$
\end{itemize}

\textbf{问题}：
\begin{enumerate}
\item 什么是阻抗和导纳？
\item 为什么运动控制器需要高阻抗（低导纳）？
\item 公式$\Delta X = Y \Delta F$如何理解？
\end{enumerate}

\section{阻抗和导纳的定义}

\subsection{基本定义}

\textbf{阻抗（Impedance）$Z$}：
\begin{equation}
Z = \frac{F}{X} \quad \text{或} \quad F = ZX
\end{equation}

\textbf{物理意义}：
\begin{itemize}
\item 阻抗表示产生单位位移所需的力
\item 高阻抗：需要很大的力才能产生小的位移
\item 低阻抗：需要小的力就能产生大的位移
\end{itemize}

\textbf{导纳（Admittance）$Y$}：
\begin{equation}
Y = \frac{X}{F} \quad \text{或} \quad X = YF
\end{equation}

\textbf{物理意义}：
\begin{itemize}
\item 导纳表示单位力产生的位移
\item 高导纳：小的力就能产生大的位移
\item 低导纳：大的力只能产生小的位移
\end{itemize}

\textbf{关系}：
\begin{equation}
Y = \frac{1}{Z} \quad \text{或} \quad Z = \frac{1}{Y}
\end{equation}

\textbf{等价表述}：
\begin{itemize}
\item 高阻抗 $\Leftrightarrow$ 低导纳
\item 低阻抗 $\Leftrightarrow$ 高导纳
\end{itemize}

\subsection{在频域中的表示}

\textbf{阻抗传递函数}：
\begin{equation}
Z(s) = \frac{F(s)}{X(s)}
\end{equation}

\textbf{导纳传递函数}：
\begin{equation}
Y(s) = \frac{X(s)}{F(s)} = \frac{1}{Z(s)}
\end{equation}

\textbf{对于线性系统}：
\begin{itemize}
\item 如果系统是线性的，阻抗和导纳是传递函数
\item 阻抗和导纳互为倒数
\end{itemize}

\section{运动控制器的目标}

\subsection{运动控制器的特点}

\textbf{目标}：
\begin{itemize}
\item 精确控制位置/运动
\item 对外力干扰不敏感
\item 即使有外力干扰，位置变化也很小
\end{itemize}

\textbf{理想情况}：
\begin{itemize}
\item 外力干扰$\Delta F$对位置的影响$\Delta X$应该尽可能小
\item 即：$\Delta X \to 0$（即使$\Delta F \neq 0$）
\end{itemize}

\subsection{公式$\Delta X = Y \Delta F$}

\textbf{公式含义}：
\begin{equation}
\Delta X = Y \Delta F
\end{equation}

其中：
\begin{itemize}
\item $\Delta F$：外力干扰（力的变化）
\item $\Delta X$：位置变化（位移的变化）
\item $Y$：导纳
\end{itemize}

\textbf{物理意义}：
\begin{itemize}
\item 外力干扰$\Delta F$通过导纳$Y$产生位置变化$\Delta X$
\item 导纳$Y$越大，相同的外力干扰$\Delta F$产生的位移变化$\Delta X$越大
\item 导纳$Y$越小，相同的外力干扰$\Delta F$产生的位移变化$\Delta X$越小
\end{itemize}

\section{为什么需要高阻抗（低导纳）？}

\subsection{数学分析}

\textbf{从公式$\Delta X = Y \Delta F$}：

如果希望位置变化$\Delta X$小，而外力干扰$\Delta F$是给定的，那么：
\begin{equation}
\Delta X = Y \Delta F \quad \Rightarrow \quad Y = \frac{\Delta X}{\Delta F}
\end{equation}

\textbf{结论}：
\begin{itemize}
\item 如果$\Delta X$要小，而$\Delta F$是固定的，那么$Y$必须小
\item 即：需要\textbf{低导纳}
\item 由于$Z = 1/Y$，低导纳对应\textbf{高阻抗}
\end{itemize}

\subsection{具体例子}

\textbf{例子1：高导纳（低阻抗）系统}

设$Y = 10$ m/N（高导纳）

如果外力干扰$\Delta F = 1$ N：
\begin{equation}
\Delta X = Y \Delta F = 10 \times 1 = 10 \text{ m}
\end{equation}

\textbf{问题}：位置变化很大（10米），这对运动控制器来说是不可接受的。

\textbf{例子2：低导纳（高阻抗）系统}

设$Y = 0.01$ m/N（低导纳）

如果外力干扰$\Delta F = 1$ N：
\begin{equation}
\Delta X = Y \Delta F = 0.01 \times 1 = 0.01 \text{ m} = 1 \text{ cm}
\end{equation}

\textbf{优点}：位置变化很小（1厘米），这对运动控制器来说是可以接受的。

\textbf{例子3：理想运动控制器}

理想情况下，$Y = 0$（零导纳，无限大阻抗）

如果外力干扰$\Delta F = 1$ N：
\begin{equation}
\Delta X = Y \Delta F = 0 \times 1 = 0 \text{ m}
\end{equation}

\textbf{优点}：位置变化为零，完全不受外力干扰影响。

\section{阻抗和导纳的物理类比}

\subsection{机械系统类比}

\textbf{弹簧系统}：
\begin{itemize}
\item \textbf{硬弹簧}（高刚度$k$）：
\begin{itemize}
\item 高阻抗：$Z = k$（大）
\item 低导纳：$Y = 1/k$（小）
\item 需要很大的力才能产生小的位移
\item 外力干扰对位置影响小
\end{itemize}

\item \textbf{软弹簧}（低刚度$k$）：
\begin{itemize}
\item 低阻抗：$Z = k$（小）
\item 高导纳：$Y = 1/k$（大）
\item 需要小的力就能产生大的位移
\item 外力干扰对位置影响大
\end{itemize}
\end{itemize}

\textbf{类比}：
\begin{itemize}
\item 好的运动控制器应该像"硬弹簧"
\item 高阻抗（低导纳），对外力干扰不敏感
\end{itemize}

\subsection{电路系统类比}

\textbf{电路系统}：
\begin{itemize}
\item \textbf{高阻抗电路}：
\begin{itemize}
\item 电流变化对电压影响小
\item 类似：力变化对位置影响小
\end{itemize}

\item \textbf{低阻抗电路}：
\begin{itemize}
\item 电流变化对电压影响大
\item 类似：力变化对位置影响大
\end{itemize}
\end{itemize}

\section{与力控制器的对比}

\subsection{力控制器的特点}

\textbf{目标}：
\begin{itemize}
\item 精确控制力
\item 对运动干扰不敏感
\item 即使有运动干扰，力变化也很小
\end{itemize}

\textbf{理想情况}：
\begin{itemize}
\item 运动干扰$\Delta X$对力的影响$\Delta F$应该尽可能小
\item 即：$\Delta F \to 0$（即使$\Delta X \neq 0$）
\end{itemize}

\textbf{公式}：
\begin{equation}
\Delta F = Z \Delta X
\end{equation}

\textbf{结论}：
\begin{itemize}
\item 如果$\Delta F$要小，而$\Delta X$是给定的，那么$Z$必须小
\item 即：需要\textbf{低阻抗}（高导纳）
\end{itemize}

\subsection{对比总结}

\begin{center}
\begin{tabular}{|l|l|l|}
\hline
\textbf{控制器类型} & \textbf{需要} & \textbf{原因} \\
\hline
运动控制器 & 高阻抗（低导纳） & $\Delta X = Y \Delta F$，$Y$小则$\Delta X$小 \\
\hline
力控制器 & 低阻抗（高导纳） & $\Delta F = Z \Delta X$，$Z$小则$\Delta F$小 \\
\hline
\end{tabular}
\end{center}

\section{实际实现}

\subsection{如何实现高阻抗}

\textbf{方法1：高增益控制}
\begin{itemize}
\item 使用大的比例增益$K_p$
\item 大的$K_p$产生大的恢复力，抵抗外力干扰
\item 结果：高阻抗
\end{itemize}

\textbf{方法2：刚性结构}
\begin{itemize}
\item 使用刚性机械结构
\item 减少柔顺性
\item 结果：高阻抗
\end{itemize}

\textbf{方法3：快速响应}
\begin{itemize}
\item 使用快速的控制回路
\item 快速检测和补偿外力干扰
\item 结果：有效的高阻抗
\end{itemize}

\subsection{阻抗的测量}

\textbf{方法}：
\begin{itemize}
\item 施加已知的外力干扰$\Delta F$
\item 测量位置变化$\Delta X$
\item 计算导纳：$Y = \Delta X / \Delta F$
\item 计算阻抗：$Z = 1/Y = \Delta F / \Delta X$
\end{itemize}

\textbf{好的运动控制器}：
\begin{itemize}
\item $Y$应该很小（低导纳）
\item $Z$应该很大（高阻抗）
\end{itemize}

\section{数学严格推导}

\subsection{线性系统的阻抗}

\textbf{假设}：系统是线性的，可以用传递函数描述。

\textbf{位置响应}：
\begin{equation}
X(s) = Y(s)F(s)
\end{equation}

\textbf{对于阶跃力干扰}：
\begin{equation}
F(s) = \frac{\Delta F}{s}
\end{equation}

\textbf{位置变化}：
\begin{equation}
\Delta X(s) = Y(s) \frac{\Delta F}{s}
\end{equation}

\textbf{稳态值}（使用终值定理）：
\begin{equation}
\Delta X_{ss} = \lim_{s \to 0} s \Delta X(s) = \lim_{s \to 0} Y(s) \Delta F = Y(0) \Delta F
\end{equation}

\textbf{结论}：
\begin{itemize}
\item 稳态位置变化$\Delta X_{ss} = Y(0) \Delta F$
\item 如果$Y(0)$小，则$\Delta X_{ss}$小
\item 因此需要低导纳（高阻抗）
\end{itemize}

\subsection{二阶系统的例子}

\textbf{二阶系统}：
\begin{equation}
m\ddot{x} + b\dot{x} + kx = f
\end{equation}

\textbf{传递函数}：
\begin{equation}
\frac{X(s)}{F(s)} = Y(s) = \frac{1}{ms^2 + bs + k}
\end{equation}

\textbf{直流增益}（$s = 0$）：
\begin{equation}
Y(0) = \frac{1}{k}
\end{equation}

\textbf{结论}：
\begin{itemize}
\item 导纳$Y(0) = 1/k$
\item 如果刚度$k$大，则导纳$Y(0)$小（低导纳）
\item 如果刚度$k$大，则阻抗$Z = k$大（高阻抗）
\item 因此，高刚度对应高阻抗（低导纳）
\end{itemize}

\section{总结}

\subsection{核心结论}

\begin{enumerate}
\item \textbf{公式}：$\Delta X = Y \Delta F$
\begin{itemize}
\item 外力干扰$\Delta F$通过导纳$Y$产生位置变化$\Delta X$
\end{itemize}

\item \textbf{运动控制器的目标}：
\begin{itemize}
\item 位置变化$\Delta X$应该尽可能小
\item 即使有外力干扰$\Delta F$
\end{itemize}

\item \textbf{需要的特性}：
\begin{itemize}
\item 低导纳$Y$（小）
\item 高阻抗$Z = 1/Y$（大）
\end{itemize}

\item \textbf{物理意义}：
\begin{itemize}
\item 高阻抗：需要很大的力才能产生小的位移
\item 低导纳：小的力只能产生小的位移
\item 结果：外力干扰对位置影响小
\end{itemize}
\end{enumerate}

\subsection{关键公式}

\begin{enumerate}
\item \textbf{导纳}：$Y = \frac{\Delta X}{\Delta F}$
\item \textbf{阻抗}：$Z = \frac{\Delta F}{\Delta X} = \frac{1}{Y}$
\item \textbf{关系}：$\Delta X = Y \Delta F$
\item \textbf{运动控制器}：需要$Y$小（低导纳），$Z$大（高阻抗）
\item \textbf{力控制器}：需要$Z$小（低阻抗），$Y$大（高导纳）
\end{enumerate}

\subsection{实际意义}

\begin{itemize}
\item \textbf{设计运动控制器时}：应该追求高阻抗（低导纳）
\item \textbf{设计力控制器时}：应该追求低阻抗（高导纳）
\item \textbf{测量系统特性时}：可以通过施加力干扰测量位置变化来计算阻抗/导纳
\item \textbf{评估控制器性能时}：阻抗/导纳是重要的性能指标
\end{itemize}

\end{document}

